% LaTeX file \f\or resume 
% This file uses the resume document class (res.cls)

\documentclass[margin]{res} 
\usepackage{color}
\usepackage{hyperref}
\usepackage{dirtytalk}
\pagestyle{plain}

% the margin option causes section titles to appear to the left of body text 
\textwidth=5.2in % increase textwidth to get smaller right margin
%\usepackage{helvetica} % uses helvetica postscript font (download helvetica.sty)
%\usepackage{newcent}   % uses new century schoolbook postscript font 

\begin{document} 
 
\name{\huge{Shoaib Jameel}\\[18pt]} % the \\[12pt] adds a blank line after name
 
\address{{\bf CONTACT INFORMATION} \\Room: 32.3035,\\Electronics and Computer Science,\\
University of Southampton,\\
Southampton,\\
Hampshire.\\
Postcode: SO17 1BJ,\\
United Kingdom\\}
\address{{}\\
 \textit{Mobile}: +44 7440532550 \\ \textit{E-mail:} M.S.Jameel@southampton.ac.uk\\
\textit{Website}: https://bashthebuilder.github.io/\\
\textit{Google Scholar}: https://scholar.google.co.uk\\/citations?user=gyy0YOIAAAAJ\&hl=en}
 
\begin{resume} 
 
\section{RESEARCH INTERESTS} 
Probabilistic Graphical Models, Approximate Posterior Inference, Bayesian Nonparametric Statistics, Probabilistic Topic Modelling, Web Search and Information Retrieval, Natural Language Processing, Social Networks and Urban Human Computing, Learning-to-Rank, Word Embeddings, Deep Neural Networks, Multi-modal Machine Learning

\section{WORK EXPERIENCE}
{\bf University of Southampton, United Kingdom} \hfill \textit{(August 2024 - Present)}\\
Electronics and Computer Science\\
Associate Professor (Teaching and Research)

{\bf University of Southampton, United Kingdom} \hfill \textit{(May 2022 - July 2024)}\\
Electronics and Computer Science\\
Lecturer (Teaching and Research)\\
\textcolor{magenta}{Finalist ``Vice Chancellor's Awards 2024 (Teaching)''}

{\bf Adani Institute for Education and Research -- AIER (Adani University), India} \hfill \textit{(June 2022 - December 2024)}\\
Invited Member, Board of Studies for Computer Science and Engineering (CSE)

{\bf University of Essex, United Kingdom} \hfill \textit{(May 2022 - December 2024)}\\
School of Computer Science and Electronic Engineering\\
Visiting Fellow

{\bf University of Essex, United Kingdom} \hfill \textit{(January 2020 - April 2022)}\\
School of Computer Science and Electronic Engineering\\
Senior Lecturer in Computer Science and Artificial Intelligence (Teaching and Research) \\
\textcolor{magenta}{Awarded the ``Outstanding Early-Career Researcher 2021'' award}

{\bf University of Kent, United Kingdom} \hfill \textit{(December 2017 - December 2019)}\\
School of Computing, Medway\\
Lecturer in Computing (Teaching and Research)\\
Served as a Deputy Admissions Officer (2018--2019) (Medway)

{\bf Cardiff University, United Kingdom} \hfill \textit{(July 2015 - November 2017)}\\
School of Computer Science and Informatics\\
Research Associate (postdoctoral researcher) with Professor Steven Schockaert

{\bf The Chinese University of Hong Kong (CUHK)} \hfill \textit{(November 2014 - July 2015)}\\
Department of Systems Engineering and Engineering Management (SEEM)\\
Postdoctoral Fellow with Prof. Wai Lam at CUHK and research collaboration with Dr Xing Xie at Microsoft Research Asia (MSRA)\\
Postdoctoral work was partly funded by Microsoft Research Asia (MSRA)

{\bf The Chinese University of Hong Kong (CUHK)} \hfill \textit{(August 2014 - October 2014)}\\
Department of Systems Engineering and Engineering Management (SEEM)\\
Research Assistant with Prof. Wai Lam and Prof. MENG, Mei Ling (Chairman and Professor, SEEM)

\section{EDUCATION}
 {\bf The Chinese University of Hong Kong (CUHK)} \hfill \textit{(August 2009 - July 2014)}\\
 Department of Systems Engineering and Engineering Management (SEEM)\\
PhD in Information Systems\\
\textbf{Thesis Title:} Latent Probabilistic Topic Discovery for Text Documents Incorporating Segment Structure and Word Order
\begin{itemize} \itemsep -13pt
\item {Advisor: Prof. Wai Lam}\\
\item {Areas of research: Web search and Information Retrieval, Topic Modeling}\\
\item {Committee: Prof. MENG, Mei Ling (Former Chairman and Professor, SEEM),  Professor CHENG, Chun Hung (Former Associate Professor, SEEM), Dr Michael C.L. Chau (Associate Professor, School of Business, Faculty of Business and Economics, The University of Hong Kong)}\\
\item Thesis was finalist in the ``Best Research Awards'' competition in entire Hong Kong
\end{itemize}

 
 {\bf Sikkim Manipal Institute of Technology (SMIT)}  \hfill  \textit{(September 2004-July 2008)}\\
Bachelor of Technology (B. Tech.) in Computer Science and Engineering (CSE)\\
Country: India
\begin{itemize} \itemsep -13pt %reduce space between items
\item Thesis Focus: Load Balancing and Scheduling in Cluster Computing\\
\item Advisors: Prof. Mrinal Kanti Ghose (Head, CSE, SMIT), Fredi B. Zarolia (Information Technology Services, Tata Steel Limited, India) \\
\item Co-Advisor: Dr Marimuthu Muruganant (\textit{Cantab.}), Research and Development Division, Tata Steel (currently Director at Steel Academics, Research and Consultancy Centre (SARCC), Sydney, Australia)\\
\item Other research topics studied: Natural Language Processing and Web search, Artificial Immune System \\
\item Ranked among the top 10\% in the entire batch \\
\item Selected based on ``merit'' through the ``All India Engineering Entrance Examination (AIEEE)'' ranking\\
\item My undergraduate thesis was awarded an \textbf{S} (\textit{Superb!}) grade
\end{itemize}

{\bf Kerala Samajam Model School (KSMS)} \hfill \textit{(March 1989 - March 2003)}\\
Country: India
\begin{itemize}
\item Indian Certificate of Secondary Education (ICSE -- 2001)
\begin{itemize}
\item Percentage - 88.88\%
\item Ranked Eighth in the entire school (\textit{out of approximately 250 students})
\end{itemize}
\item Indian School Certificate (ISC -- 2003)
\begin{itemize}
\item Percentage - 82.00\%
\item Ranked Fourth in the Science stream (\textit{out of approximately 120 students}). \newline \textit{(Missed out on the third spot by a score of 0.5)}
\end{itemize}
\end{itemize}

\section{CONFERENCE PUBLICATIONS}
\begin{enumerate} \itemsep -10pt

\item Ali Almutairi, Aditya Joshi, Gelareh Mohammadi, Abdullah Alsuhaibani, \textbf{Shoaib Jameel}, Imran Razzak. \textit{FLICK: Few-Label Incremental Learning for Low-Resource Dialects}. The International Conference Language Technologies for Low-resource Languages (LaTeLL 2026). \\

\item Ghadir Alselwi, Hao Xue, \textbf{Shoaib Jameel}, Basem Suleiman, Flora D. Salim, Imran Razzak. 2026. \textit{Long Context Modeling with Ranked Memory-Augmented Retrieval}. The 64th Annual Meeting of the Association for Computational Linguistics (ACL) full paper. \\

\item Jianghao Wu, Feilong Tang, Yulong Li, Ming Hu, Haochen Xue, \textbf{Shoaib Jameel}, Zongyuan Ge, Yutong Xie, Imran Razzak. \textit{TAGS: A Test-Time Generalist–Specialist Framework with Retrieval-Augmented Reasoning and Verification}. 2026. The 64th Annual Meeting of the Association for Computational Linguistics (ACL) full paper. \\

\item Ubaid Azam, Imran Razzak, \textbf{Shoaib Jameel}. 2026. \textit{Lightweight Cross-Lingual Federated Prompt Tuning for Low-Resource Languages}. Language Resources and Evaluation Conference 2026. \\

\item Youssef Mohamed, Noran Mohamed, Khaled Abouhashad, Feilong Tang, Sara Atito, \textbf{Shoaib Jameel}, Imran Razzak, and Ahmed B. Zaky. 2025. \textit{DeepChest: Dynamic Gradient-Free Task Weighting for Effective Multi-Task Learning in Chest X-ray Classification}. IEEE BigData 2025 Conference. \\

\item Inamullah, Imran Razzak, \textbf{Shoaib Jameel}. 2025. \textit{Retinal–Lipidomics Associations as Candidate Biomarkers for Cardiovascular Health}. 6th International Conference on Medical Imaging and Computer-Aided Diagnosis (MICAD 2025). \\

\item Lingzhi Shen, Xiaohao Cai, Yunfei Long, Imran Razzak, Guanming Chen, \textbf{Shoaib Jameel}. 2025. \textit{CALM: Culturally Self-Aware Language Models}. The Thirty-Ninth Annual Conference on Neural Information Processing Systems (NeurIPS 2025). \\

\item Guanming Chen, Lingzhi Shen, Xiaohao Cai, Imran Razzak, \textbf{Shoaib Jameel}. 2025. \textit{HIPPD: Brain-Inspired Hierarchical Information Processing for Personality Detection}. The 17th Asian Conference on Machine Learning (ACML--2025). \\

\item  Lingzhi Shen, Xiaohao Cai, Yunfei Long, Imran Razzak, Guanming Chen and \textbf{Shoaib Jameel}. 2025. \textit{EmoPerso: Enhancing Personality Detection with Self-Supervised Emotion-Aware Modelling}. The Conference on Information and Knowledge Management (CIKM -- 2025). \\

\item Ant\'onio Correia, Tommi Karkkainen, \textbf{Shoaib Jameel}, Daniel Schneider, Pedro Antunes, Benjamim Fonseca, Andrea Grover. 2025. \textit{A Pipeline for AI-Based Quantitative Studies of Science Enhanced by Crowdsourced Inferential Modelling}. 23rd International Conference on Hybrid Intelligent Systems (HIS 2025), December 11-13, 2023, Volume 4: Machine Learning Applications. \\

\item Abdulsalam Alharbi, \textbf{Shoaib Jameel}, Basem Suleiman and Imran Razzak. 2025. \textit{Enhancing Large Language Models for Arabic Dialects Using Knowledge-Based Rethinking and Contrastive Learning}. 17th International Conference on Social Networks Analysis and Mining (ASONAM 2025). \\

\item Mozhgan Talebpour, Yunfei Long, Alba	Garcia Seco De Herrera, \textbf{Shoaib Jameel}. 2025. \textit{Bias in Language Models: Interplay of Architecture and Data?} The 48th International ACM SIGIR Conference on Research and Development in Information Retrieval. \\

\item Lingzhi Shen, Yunfei Long, Xiaohao Cai, Guanming Chen, Imran Razzak, and \textbf{Shoaib Jameel}. 2025. \textit{Less but Better: Parameter-Efficient Fine-Tuning of Large Language Models for Personality Detection}. International Joint Conference on Neural Networks (IJCNN). Accepted Full paper. \\

\item Lingzhi Shen, Yunfei Long, Xiaohao Cai, Guanming Chen, Yuhan Wang, Imran Razzak, \textbf{Shoaib Jameel}. 2025. \textit{Self-Supervised Dynamic Optimization for Graph-Based Personality Detection}. IEEE International Conference on Multimedia\&Expo 2025. Accepted Full paper. \\

\item Usman Naseem, Qi Zhang, Shoujin Wang, Liang Hu and \textbf{Shoaib Jameel}. 2025. \textit{DiGrI: Distorted Greedy Approach for Human-Assisted Online Suicide Ideation Detection.} The Web Conference. \\

\item Shijing Chen, Mohamed Reda Bouadjenek, Usman Naseem, Basem Suleiman, \textbf{Shoaib Jameel}, Flora Salim, Hakim Hacid and Imran Razzak. \textit{Leveraging Taxonomy and LLMs for Improved Multimodal Hierarchical Classification.} The 31st International Conference on Computational Linguistics (COLING 2025). \newline

\item Ubaid Azam, Imran Razzak, Shelly Vishwakarma and \textbf{Shoaib Jameel}. 2025. \textit{Uncertainty Modelling in Under-Represented Languages with Bayesian Deep Gaussian Processess}. The 31st International Conference on Computational Linguistics (COLING 2025). \newline

\item Ubaid Azam, Imran Razzak, Shelly Vishwakarma, Hakim Hacid, Dell Zhang, \textbf{Shoaib Jameel}. 2025. \textit{BAKER: Bayesian Kernel Uncertainty in Domain-Specific Document Modelling}. The Eighteenth International Conference on Web Search and Data Mining (WSDM 2025). Accepted. Full paper. \textcolor{magenta}{In collaboration with Thomson Reuters Labs} \newline

\item Lingzhi Shen, Yunfei Long, Xiaohao Cai, Imran Razzak, Guanming Chen, Kang Liu, \textbf{Shoaib Jameel}. 2025. \textit{GAMED: Knowledge Adaptive Multi-Experts Decoupling for Multimodal Fake News Detection}.  The Eighteenth International Conference on Web Search and Data Mining (WSDM 2025). Accepted. Full paper. \newline

\item Ubaid Azam, Imran Razzak, Shelly Vishwakarma, Hakim Hacid, Dell Zhang and \textbf{Shoaib Jameel}. 2024. \textit{Would You Trust an AI Doctor? Building Reliable Medical Predictions with Kernel Dropout Uncertainty}. The 25th International Web Information Systems Engineering Conference (WISE 2024). Accepted. Full paper. \textcolor{magenta}{In collaboration with Thomson Reuters Labs} \newline

\item Abdullah Alsuhaibani, Imran Razzak, \textbf{Shoaib Jameel}, Xianzhi Wang and Guandong Xu. 2024. \textit{CLIMB: Imbalanced Data Modelling using Contrastive Learning with Limited Labels}. The 25th International Web Information Systems Engineering Conference (WISE 2024). Accepted. Full paper. \newline

\item Ant\'onio Correia, Tommi K\"arkk\"ainen, \textbf{Shoaib Jameel}, Daniel Schneider, Pedro Antunes, Benjamim Fonseca, Andrea Grover. \textit{A DSR Approach for the Design of a Hybrid Crowd-Algorithmic RLHF-based Scientometric System}. The 23rd International Conference on Hybrid Intelligent Systems (HIS 2023). Accepted. Full paper. \newline

\item Abdullah Alsuhaibani, Hamad Zogan, Imran Razzak, \textbf{Shoaib Jameel}, Guandong Xu. 2023. \textit{IDoFew: Intermediate Training Using Dual-Clustering in Language Models for Few Labels Text Classification}. The 17th ACM International Conference Web Search and Data Mining (WSDM 2023). Accepted. Full paper. \newline

\item Ant\'onio Correia, Andrea Grover, \textbf{Shoaib Jameel}, Benjamim Fonseca. \textit{A DSR Approach for the Design of a Hybrid Crowd-Algorithmic RLHF-based Scientometric Data Analysis System}. 2023. International Conference on Industry Sciences and Computer Science Innovation. Accepted. Full paper. \newline

\item Mozhgan Talebpour, Alba Garcia Seco de Herrera, \textbf{Shoaib Jameel}. 2023. \textit{Topics in Contextualised Attention Embeddings}. The 45th European Conference on Information Retrieval (ECIR). Accepted. Full paper. \newline

\item Kun Yan, Chenbin Zhang, Jun Hou, Ping Wang, Zied Bouraoui, \textbf{Shoaib Jameel}, Steven Schockaert. 2022. \textit{Inferring Prototypes for Multi-Label Few-Shot Image Classification with Word Vector Guided Attention}. Thirty-Sixth AAAI Conference on Artificial Intelligence, Vancouver. Accepted. Full paper. \newline

\item António Correia, Benjamim Fonseca, Hugo Paredes, Ramon Chaves, Daniel Schneider, and \textbf{Shoaib Jameel}. 2021. \textit{Determinants and Predictors of Intentionality and Perceived Reliability in Human-AI Interaction as a Means for Innovative Scientific Discovery}. The 3rd International Workshop on Big Data Tools, Methods, and Use Cases for Innovative Scientific Discovery (BTSD) 2021. Accepted. Short paper. \newline

\item Li He, Hongxu Chen, Dingxian Wang, \textbf{Shoaib Jameel}, Philip Yu and Guandong Xu. 2021. \textit{Click-Through Rate Prediction with Multi-Modal Hypergraphs}. Conference on Information and Knowledge Management (CIKM 2021). Accepted. Full paper. (\textcolor{magenta}{Nominated for the ``Best Paper Award'', ``Winner in the ACS Digital Disruptors Awards 2022'', ``Selected by the eBay leadership team as a Top 3 Leaders Choice Award recipient out of 2000 international submissions.'', ``Featured in the \href{https://www.bandt.com.au/reaching-gen-z-a-new-tool-for-marketers/}{media}''}). \textit{CIKM is a major/premier conference in data mining and knowledge management.} \newline

\item Elena Barry, \textbf{Shoaib Jameel}, Haider Raza. 2021. \textit{Emojional: Emoji Embeddings}. 20th Annual UK Workshop on Computational Intelligence (UKCI 2021). Accepted. Full paper. \textit{[Authored by a final-year undergraduate student. We were recently contacted by \textcolor{magenta}{BBC Radio 4} to talk about this research.]} \newline

\item Hamad Zogan, Imran Razzak, \textbf{Shoaib Jameel} and Guandong Xu. 2021. \textit{DepressionNet: Learning Multi-modalities with User Post Summarization for Depression Detection on Social Media}. The 44th International ACM SIGIR Conference on Research and Development in Information Retrieval (SIGIR 2021). Accepted. Full paper. \newline

\item Kun Yan, Zied Bouraoui, Ping Wang, \textbf{Shoaib Jameel}, Steven Schockaert. 2021. \textit{Aligning Visual Prototypes with BERT Embeddings for Few-Shot Learning}. ACM International Conference on Multimedia Retrieval 2021 (ICMR 2021). Accepted. Full paper. \newline

\item Ant\'onio Correia, Daniel Schneider, \textbf{Shoaib Jameel}, Hugo Paredes, and Benjamim Fonseca. 2021. \textit{Experimental Investigation of the Factors Influencing Researchers' Adoption of Crowdsourcing and Machine Learning: An Initial Survey}. 20th International Conference on Intelligent Systems Design and Applications (ISDA 2020). Online conference. Full paper. \newline

\item Ant\'onio Correia, Diogo Guimaraes, Dennis Paulino, \textbf{Shoaib Jameel}, Daniel Schneider, Benjamim Fonseca and Hugo Paredes. 2021. \textit{AuthCrowd: Author Name Disambiguation and Entity Matching using Crowdsourcing}. 24th IEEE International Conference on Computer Supported Cooperative Work in Design (IEEE CSCWD 2021). Online conference. Full paper. \newline

\item Kun Yan, Zied Bouraoui, Ping Wang, \textbf{Shoaib Jameel}, Steven Schockaert. 2021. \textit{Few-Shot Image Classification With Multi-Facet Prototypes}. International Conference on Acoustics, Speech, \& Signal Processing (ICASSP). Online conference. Full paper. \newline

\item Ant\'onio Correia, \textbf{Shoaib Jameel}, Daniel Schneider, Hugo Paredes, and Benjamim Fonseca. 2020. \textit{A Workflow-Based Methodological Framework for Hybrid Human-AI Enabled Scientometrics}. IEEE BigData. 2021. Online conference. Full paper. \newline

\item Zihao Fu, Lidong Bing, Wai Lam, \textbf{Shoaib Jameel}. \textit{Dynamic Topic Tracker for KB-to-Text Generation}. 2020. International Conference on Computational Linguistics (COLING). Online conference. Full paper. \newline

\item Ant\'onio Correia, Benjamim Fonseca, Hugo Paredes, Daniel Schneider, \textbf{Shoaib Jameel}. \textit{On the Socio-Technical Aspects of Crowdsourcing as a Mechanism for Academic Study: Conceptual Framework and Scoping Review}. 2020. The Hawaii International Conference on System Sciences (HICSS), Hawaii, USA. (\textcolor{magenta}{Nominated for the ``Best Paper Award''}) \newline

\item Ant\'onio Correia, Dennis Paulino, Hugo Paredes, Benjamim Fonseca, \textbf{Shoaib Jameel}, Daniel Schneider and Jano de Souza. \textit{Scientometric Research Assessment of IEEE CSCWD Conference Proceedings: An Exploratory Analysis from 2001 to 2019}. Accepted. Full paper. Dalian, China. \textit{Conference postponed due to SARS-COV-2 outbreak.} \newline

\item Ant\'onio Correia, Hugo Paredes, Benjamim Fonseca, \textbf{Shoaib Jameel}. \textit{SciCrowd: A crowd-AI system for supporting crowd work in scientific research}. 2019. In: Science and Technology Summit in Portugal. Poster. \newline

\item Ant\'onio Correia, Benjamim Fonseca, Hugo Paredes, Daniel Schneider, \textbf{Shoaib Jameel}. \textit{Development of a Crowd-Powered System Architecture for Knowledge Discovery in Scientific Domains}. 2019. 2019 IEEE International Conference on Systems, Man, and Cybernetics (SMC), pp. 1372--1377, Bari, Italy. \newline

\item Ant\'onio Correia, Hugo Paredes, Daniel Schneider, \textbf{Shoaib Jameel}, Benjamim Fonseca. \textit{Towards Hybrid Crowd-AI Centered Systems: Developing an Integrated Framework from an Empirical Perspective}. 2019. 2019 IEEE International Conference on Systems, Man, and Cybernetics (SMC), pp. 4013--4018, Bari, Italy. \newline

\item \textbf{Shoaib Jameel}, and Steven Schockaert. 2019. \textit{Word and Document Embedding with vMF-Mixture Priors on Context Word Vectors}.  In the 57th Annual Meeting of the Association for Computational Linguistics (ACL), pp. 3319--3328, Florence, Italy. \newline

\item Jose Camacho-Collados, Luis Espinosa-Anke, \textbf{Shoaib Jameel}, Steven Schockaert. \textit{A Latent Variable Model for Learning Distributional Relation Vectors}. 2019. In the 28th International Joint Conference on Artificial Intelligence (IJCAI), 4911--4917, Macao, China, \newline

\item Ant\'onio Correia, \textbf{Shoaib Jameel}, Daniel Schneider, Benjamim Fonseca and Hugo Paredes. 2019. \textit{The Effect of Scientific Collaboration on CSCW Research: A Scientometric Study}. In the 23rd IEEE International Conference on Computer Supported Cooperative Work in Design (IEEE CSCWD), pp. 149--161, Porto, Portugal. \newline

\item \textbf{Shoaib Jameel}, Zihao Fu, Bei Shi, Wai Lam, and Steven Schockaert. 2019. \textit{Word Embedding as Maximum A Posteriori Estimation}. In the 33rd AAAI Conference on Artificial Intelligence (AAAI), pp. 6562--6569, Honolulu, Hawaii, USA. [\href{https://kar.kent.ac.uk/70009/7/AAAI-2019-Camera-Ready\%20\%281\%29.pdf}{PDF}] \newline

\item Zied Bouraoui, \textbf{Shoaib Jameel}, and Steven Schockaert. 2018. \textit{Relation Induction in Word Embeddings Revisited}. In the 27th International Conference on Computational Linguistics (COLING), pp. 1627--1637. Santa Fe, New Mexico. [\href{https://arxiv.org/abs/1708.06266}{PDF}] (\textcolor{magenta}{One of the ``Area Chair Favourites''}) \newline

\item \textbf{Shoaib Jameel}, Zied Bouraoui and Steven Schockaert. 2018. \textit{Unsupervised Learning of Distributional Relation Vectors}. In the 56th Annual Meeting of the Association for Computational Linguistics  (ACL), pp. 23--33. Melbourne, Australia. [\href{https://drive.google.com/open?id=1zwjJFiu6yaq-yq2as9liFkTagYlSl8x9}{PDF}] \newline

\item \textbf{Shoaib Jameel}, and Steven Schockaert. 2017. \textit{Modeling Context Words as Regions: An Ordinal Regression Approach to Word Embedding}.  In the SIGNLL Conference on Computational Natural Language Learning (CoNLL), pp. 123--133, Vancouver, Canada. [\href{https://drive.google.com/open?id=18WMMGB-ohy_EkA_nyl3n8wDWMcLNX7HH}{PDF}] \newline

\item  Bei Shi, Wai Lam, \textbf{Shoaib Jameel}, Steven Schockaert and Kwun Ping Lai. 2017. \textit{Jointly Learning Word Embeddings and Latent Topics}. In the 40th International ACM SIGIR Conference on Research and Development in Information Retrieval (SIGIR), pp. 375--384, Tokyo, Japan. [\href{https://drive.google.com/open?id=1I_fmdLmae0SE1HEzBRO7Nv-RedkNlnAY}{PDF}] \newline

\item \textbf{Shoaib Jameel}, Zied Bouraoui, and Steven Schockaert. 2017. \textit{MEmbER: Max-Margin Based Embeddings for Entity Retrieval}. In the 40th International ACM SIGIR Conference on Research and Development in Information Retrieval (SIGIR), pp. 783--792, Tokyo, Japan. [\href{https://drive.google.com/open?id=1QVXU8E6VtG8TTKbfTeU4DvCifh8keWQa}{PDF}] \newline

\item Zied Bouraoui, \textbf{Shoaib Jameel}, and Steven Schockaert. 2017. \textit{Inductive Reasoning about Ontologies Using Conceptual Spaces}. In Proceedings of the 31st AAAI Conference on Artificial Intelligence (AAAI), pp. 4364--4370, San Francisco, California, USA. [\href{https://drive.google.com/open?id=1CDoH5kqyN3UnqDgwSGaGqBOxnT7vvime}{PDF}] \newline

\item \textbf{Shoaib Jameel}, Steven Schockaert. 2016. \textit{\uppercase{D}-\uppercase{G}loVe: \uppercase{A} Feasible Least Squares Model for Estimating Word Embedding Densities}. In Proceedings of the 26th International Conference on Computational Linguistics (COLING), pp. 1849--1860, Osaka, Japan. [\href{https://drive.google.com/open?id=1kgbHGwtLJBZPXKQWAMXp7WzNSLX9giFL}{PDF}] \newline

\item \textbf{Shoaib Jameel}, Yi Liao, Wai Lam, Steven Schockaert, Xing Xie. 2016. \textit{Exploring Urban Lifestyles Using a Nonparametric Temporal Graphical Model}. In Proceedings of the ACM SIGIR International Conference on the Theory of Information Retrieval (ICTIR), pp. 251--260, Newark, Delaware, USA. [\href{https://drive.google.com/open?id=1MrnfHbZO9jGJao9diVJihHmKIDj9KXRp}{PDF}] \newline

\item Yi Liao, Wai Lam, \textbf{Shoaib Jameel}, Steven Schockaert, Xing Xie. 2016. \textit{Who Wants to Join Me? Companion Recommendation in Location-Based Social Networks}. In Proceedings of the ACM SIGIR International Conference on the Theory of Information Retrieval (ICTIR), pp. 271--280, Newark, Delaware, USA. [\href{https://drive.google.com/open?id=1iPLZ8EDQaQYkjpZ6FTlK9ZPksf_Ki8wN}{PDF}] \newline

\item \textbf{Shoaib Jameel}, and Steven Schockaert. 2016. \textit{Entity Embeddings with Conceptual Subspaces as a Basis for Plausible Reasoning}. In Proceedings of the European Conference on Artificial Intelligence (ECAI), pp. 1353--1361, The Hague, Netherlands. [\href{https://arxiv.org/abs/1602.05765}{PDF}] \newline

\item Steven Schockaert, and \textbf{Shoaib Jameel}. 2016. \textit{Plausible Reasoning based on Qualitative Entity Embeddings}. In Proceedings of the 25th International Joint Conference on Artificial Intelligence (IJCAI). (\textit{Invited Paper}), pp. 4078--4081, New York, USA. [\href{https://www.ijcai.org/Proceedings/16/Papers/605.pdf}{PDF}] \newline

\item Pengfei Liu, \textbf{Shoaib Jameel}, King Keung Wu, and Helen Meng. 2016. \textit{Learning Track Representation and Trends for Conference Analytics}. In Proceedings of the Hawaii International Conference on System Sciences (HICSS-49), pp. 1671--1680, Hawaii, USA. [\href{https://drive.google.com/file/d/1mRGqiIsAY1hFgvnr5zSQ_m5c0FTgDuZQ/view?usp=sharing}{PDF}]  \newline

\item Yi Liao, \textbf{Shoaib Jameel}, Wai Lam, and Xing Xie. 2015. \textit{Abstract Venue Concept Detection from Location-Based Social Networks}. In Proceedings of the Asia Information Retrieval Societies Conference (AIRS), pp. 147--157, Brisbane, Australia. [\href{https://drive.google.com/file/d/1aRgJWZJPdfd_cMPB7IucGDO9FQIJ-SFG/view?usp=sharing}{PDF}]  \newline

\item \textbf{Shoaib Jameel}, Wai Lam, Steven Schockaert, and Lidong Bing. 2015. \textit{A Unified Posterior Regularized Topic Model with Maximum Margin for Learning-to-Rank}. In Proceedings of the 24th ACM International Conference on Information and Knowledge Management (CIKM), pp. 103--112, Melbourne, Australia. [\href{https://drive.google.com/file/d/1FlUuVeUd0eSb7aNrdJZgm5LdRKLvNDDT/view?usp=sharing}{PDF}]  \newline

\item Pengfei Liu, \textbf{Shoaib Jameel}, Wai Lam, Bin Ma, and Helen Meng. 2015. \textit{Topic Modeling for Conference Analytics}. In Proceedings of the 16th Annual Conference of the International Speech Communication Association (INTERSPEECH). [\href{https://drive.google.com/file/d/11Bz2_A6kpyR01ED8rnfECo4aDNKvH-yF/view?usp=sharing}{PDF}]  \newline

\item \textbf{Shoaib Jameel}, Wai Lam, and Lidong Bing. 2015. \textit{Nonparametric Topic Modeling using Chinese Restaurant Franchise with Buddy Customers}. In Proceedings of the 37th European Conference on Information Retrieval (ECIR), pp. 648--659, Vienna, Austria. [\href{https://drive.google.com/file/d/1m52uOZw3tPxPp3LkPfqEyR-LkPTi_m5q/view?usp=sharing}{PDF}]  \newline

\item Lidong Bing, Wai Lam, \textbf{Shoaib Jameel}, and Chunliang Lu. 2014. \textit{Website Community Mining from Query Logs with Two-phase Clustering}. In Proceedings of the 15th International
Conference on Intelligent Text Processing and Computational Linguistics (CICLing), pp. 201--212, Kathmandu, Nepal. [\href{https://drive.google.com/file/d/1INsfUIRtjIuCzLXa_dlmkR_B07lp2tC9/view?usp=sharing}{PDF}]  \newline

\item  \textbf{Shoaib Jameel}, and Wai Lam. 2013. \textit{A Nonparametric N-Gram Topic Model with Interpretable Latent Topics}. In Proceedings of the Ninth Asia Information Retrieval Societies Conference (AIRS), pp. 74--85, Singapore. [\href{https://drive.google.com/file/d/18jXDQmbGjBDzrJXSJ9tRLcRa0DH_MLVJ/view?usp=sharing}{PDF}]  \newline

\item \textbf{Shoaib Jameel}, and Wai Lam. 2013. \textit{An Unsupervised Topic Segmentation Model Incorporating Word Order}. In Proceedings of the 36th Special Interest Group on Information Retrieval (ACM-SIGIR), pp. 203--212, Dublin, Ireland. [\href{https://drive.google.com/file/d/18F6SyQP5ipi4MUbXhqAlsi92Ri7T3d3K/view?usp=sharing}{PDF}] \\

\item  \textbf{Shoaib Jameel}, and Wai Lam. 2013. \textit{An N-gram Topic Model for Time-Stamped Documents}. In Proceedings of the 35th European Conference on Information Retrieval (ECIR), pp. 292--304, Moscow, Russia. [\href{https://drive.google.com/file/d/1wR1Yeyw5l9rf1UD1SxpWTfUaUo5P34Eo/view?usp=sharing}{PDF}]\\

\item  \textbf{Shoaib Jameel}, Xiaojun Qian, and Wai Lam. 2012. \textit{N-gram Fragment Sequence Based Unsupervised Domain-Specific Document Readability}. In Proceedings of the 24th International Conference on Computational Linguistics (COLING), pp. 1309--1326, Mumbai, India. [\href{http://www.aclweb.org/anthology/C12-1080}{PDF}]\\

\item  \textbf{Shoaib Jameel}, Wai Lam, Xiaojun Qian, and Ching-man Au Yeung. 2012. \textit{An unsupervised technical difficulty ranking model based on conceptual terrain in the latent space}. In Proceedings of the 12th ACM/IEEE-CS Joint Conference on Digital Libraries (JCDL), pp. 351--352, Washington D.C., USA. [\href{https://drive.google.com/file/d/1uLTRvw9sQchIj0FF3A2tWTduNGedOIo4/view?usp=sharing}{PDF}]\\

\item  \textbf{Shoaib Jameel}, and Xiaojun Qian. 2012. \textit{An Unsupervised Technical Readability Ranking Model by Building a Conceptual Terrain in LSI}. In Proceedings of the International Conference on Semantics, Knowledge and Grids (SKG), pp. 39--46, Beijing, China. [\href{https://drive.google.com/file/d/1z-6qqTff1fo8LbmOfcTyIps6WGGNizd1/view?usp=sharing}{PDF}]\\

\item  \textbf{Shoaib Jameel}, Wai Lam, and Xiaojun Qian. 2012. \textit{Ranking Text Documents based on Conceptual Difficulty using Term Embedding and Sequential Discourse Cohesion}. In Proceedings of the 2012 IEEE/WIC/ACM International Conference on Web Intelligence (WI), pp. 145--152, Macau, China. [\href{https://drive.google.com/file/d/1eGZ6GZ8-tf2pWaQtCutobVG5kleL6yfG/view?usp=sharing}{PDF}]\\

\item  \textbf{Shoaib Jameel}, Wai Lam, Ching-man Au Yeung, and Sheaujiun Chyan. 2011. \textit{An Unsupervised Ranking Method based on a Technical Difficulty Terrain}. In Proceedings of the 20th ACM International Conference on Information and Knowledge Management (CIKM), pp. 1989--1992, Glasgow, Scotland. (\textit{\textcolor{red}{Winner: Best Paper Award at the Beijing-Hong Kong International Doctoral Forum-2011}}) \\

\item \textbf{M. Shoaib Jameel}, Anubhav, Nilesh Singh, Nitin Kumar Singh, Tejbanta Singh Chingtham, and M. K. Ghose. 2009. \textit{An Intelligent Automatic Text Summarizer}. In Proceedings of the First International Conference on Intelligent Human Computer Interaction, pp. 223--230, Hyderabad, India.\\

\item \textbf{M. Shoaib Jameel}, Amar Akshat, and Tejbanta Singh Chingtham. 2008. \textit{Enhancements in Query Evaluation and Page Summarization of The Thinking Algorithm}. In Proceedings of the International Symposium on Information Technology (ITSim), pp. 1--8, Kuala Lumpur, Malaysia.
\end{enumerate}

\section{JOURNAL PUBLICATIONS}
\begin{enumerate} \itemsep -10pt
    \item Abdullah Alsuhaibani, Abdulrahman Alalawi, Imran Razzak, \textbf{Shoaib Jameel}, Xianzhi Wang, and Guandong Xu. 2026. \textit{Cluster-SCP: Similarity and Contrastive Learning to Enhance Pseudo Labels for Fine Tuning under Few Labels}. World Wide Web Journal. \\
    
    \item Inamullah, Imran Razzak, \textbf{Shoaib Jameel}. 2025. \textit{The Eye as a Window to Systemic Health: A Survey of Retinal Imaging from Classical Techniques to Oculomics}. The Journal of Precision Medicine: Health and Disease. \textcolor{red}{Featured on: [\href{https://www.forbes.com/sites/williamhaseltine/2025/11/05/how-artificial-intelligence-makes-eye-exams-a-gateway-to-whole-body-wellness/}{Forbes}]} \\
    \item Kun Yan, Zied Bouraoui, Fangyun Wei, Chang Xu, Ping Wang, \textbf{Shoaib Jameel}, Steven Schockaert. \textit{Modelling Multi-modal Cross-interaction for Multi-label Few-shot Image Classification Based on Local Feature Selection}. 2025. Transactions on Multimedia Computing Communications and Applications. \newline
    \item Mahsa Abazari Kia, Aygul Garifullina, Mathias Kern, Jon Chamberlain, \textbf{Shoaib Jameel}. 2024. \textit{Question-Driven Text Summarization Using an Extractive-Abstractive Framework}. Computational Intelligence. Accepted. \newline
    \item Guangming Huang, Yingya Li, \textbf{Shoaib Jameel}, Yunfei Long, and Giorgos Papanastasiou. 2024. \textit{From Explainable to Interpretable Deep Learning for Natural Language Processing in Healthcare: How Far from Reality?}. Computational and Structural Biotechnology Journal. Accepted. \newline
    \item Ant\'onio Correia, Diogo Guimar\~aes, Hugo Paredes, Benjamim Fonseca, Dennis Paulino, Lu\'is Trigo, Pavel Brazdil, Daniel Schneider, Andrea Grover, and \textbf{Shoaib Jameel}. \textit{NLP-Crowdsourcing Hybrid Framework for Inter-Researcher Similarity Detection}. IEEE Transactions on Human-Machine Systems. Accepted. \newline
    \item Ant\'onio Correia, Andrea Grover, \textbf{Shoaib Jameel}, Daniel Schneider, Pedro Antunes, and Benjamim Fonseca1. 2023. \textit{A Hybrid Human-AI Tool for Scientometric Analysis}. Artificial Intelligence Review. Accepted. \newline
    \item Sujit Roy, Vishal Gaur, Haider Raza, \textbf{Shoaib Jameel}. 2023. \textit{CLEFT: Contextualised Unified Learning of User Engagement in Video Lectures with Feedback}.  IEEE Access. Accepted. \newline
    \item Hamad Zogan, Imran Razzak, \textbf{Shoaib Jameel}, Guandong Xu. 2023. \textit{Convolutional Attention Based Multimodal Network for Depression Detection on Social Media and Its Impact During Pandemic}. the IEEE Journal of Biomedical and Health Informatics (J-BHI). Accepted. \newline
    \item Li He, Guandong Xu, \textbf{Shoaib Jameel}, Xianzhi Wang, Hongxu Chen. 2022. \textit{Graph-Aware Deep Fusion Networks for Online Spam Review Detection}. IEEE Transactions on Computational Social Systems. Accepted. \newline
    \item Mahsa Abazari Kia, Aygul Garifullina, Mathias Kern, Jon Chamberlain, \textbf{Shoaib Jameel}. 2022. \textit{Adaptable Closed-Domain Question Answering Using Contextualized CNN-Attention Models and Question Expansion}. IEEE Access. Accepted with minor revisions. \newline
    \item Alberto Jose Benayas, Reyhaneh Hashempour, Damian Rumble, \textbf{Shoaib Jameel}, Renato Amorim. 2021. \textit{Unified Transformer Multi-task Learning for Intent Classification with Entity Recognition}. IEEE Access. vol. 9., pp. 147306--147314. \textcolor{magenta}{Model is currently deployed at BT (formerly British Telecom) and is being used by their development team.}  \newline
    \item Hamad Zogan, Imran Razzak, Xianzhi Wang, \textbf{Shoaib Jameel}, Guandong Xu. 2021. \textit{Explainable Depression Detection with Multi-Modalities Using a Hybrid Deep Learning Model on Social Media}. World Wide Web Journal. \\
    \item Jianlong Zhou, Hamad Zogan, Shuiqiao Yang, \textbf{Shoaib Jameel}, Guandong Xu, Fang Chen. 2021. \textit{Detecting Community Depression Dynamics Due to COVID-19 Pandemic in Australia}. IEEE Transactions on Computational Social Systems. vol. 8, no. 4, pp. 982-991, Aug. 2021, doi: 10.1109/TCSS.2020.3047604 \\
    \item Ant\'onio Correia, Daniel Schneider, \textbf{Shoaib Jameel}, Hugo Paredes, Benjamim Fonseca. 2019. \textit{`Caught Between Worlds': Leveraging Crowd Science and AI for Supporting Scientific Work Practices}. Computer Supported Cooperative Work: The Journal of Collaborative Computing and Work Practices \\

\item \textbf{Shoaib Jameel}, Wai Lam, and Lidong Bing. 2015. \textit{Supervised Topic Models with Word Order Structure for Document Classification and Retrieval Learning}. In Information Retrieval Journal (IRJ). 18 (4): 283-330.\\
\item Lidong Bing, Shan Jiang, Wai Lam, Yan Zhang, and \textbf{Shoaib Jameel}. 2015. \textit{Adaptive Concept Resolution for Document Representation and Its Applications in Text Mining}. In Knowledge-Based Systems 74 (2015): 1-13.\\
\item Lidong Bing, Wai Lam, Tak-Lam Wong, and \textbf{Shoaib Jameel}. 2015. \textit{Web Query Reformulation via Joint Modeling of Latent Topic Dependency and Term Context}. In ACM Transactions on Information Systems (TOIS). Volume 33, Issue 2. \\

\item  \textbf{M. Shoaib Jameel}, and Tejbanta Singh Chingtham. 2009. \textit{Compounded Uniqueness Level: Geo-Location Indexing Using Address Parser}. In International Journal of Computer Theory and Engineering. pp. 27--34\\

\item  \textbf{M. Shoaib Jameel}, Marimuthu Muruganant, and Tejbanta Singh Chingtham. 2009. \textit{Deploying CPU Load Balancing in the Linux Cluster Using Non-Repetitive CPU Selection}. In International Journal of Computer Theory and Electrical Engineering, pp. 228--234.
\end{enumerate}

\section{BOOK CHAPTERS}
\begin{enumerate} \itemsep -10pt
    \item Ant\'onio Correia, \textbf{Shoaib Jameel}, Hugo Paredes, Benjamim Fonseca, Daniel Schneider 2019. \textit{Hybrid Machine-Crowd Interaction for Handling Complexity in Science: A Scaffolding Design Framework}. Accepted. Springer Human-Computer Interaction Series
\end{enumerate}

\section{WORKSHOP PUBLICATIONS}
\begin{enumerate}
\item Inam Ullah, Imran Razzak, \textbf{Shoaib Jameel}. 2026. \textit{RetiSEM: Generalising Causal Models for Fragmented Biomedical Data}. Fourth International Workshop on Generalizing from Limited Resources in the Open World Workshop at International Joint Conference on Artificial Intelligence (IJCAI) 2026

\item Niranjana Arun Menon, Md Rafiqul Islam, Mohamed Reda Bouadjenek, \textbf{Shoaib Jameel}, Imran Razzak. 2026. \textit{Exploring the Influence of Lifestyle, Social Health, and Demographic Factors on Psychological Well-being and Engagement Levels}. Fourth International Workshop on Multimodal Content Analysis for Social Good (WebConf 2026).
\item Abdulsalam Obaid Alharbi, Abdullah Alsuhaibani, Abdulrahman Abdullah Alalawi, Usman Naseem, \textbf{Shoaib Jameel}, Salil Kanhere, Imran Razzak. 2025. \textit{Evaluating Large Language Models on Health-Related Claims Across Arabic Dialects}. The 1st Workshop on NLP for Languages Using Arabic Script at The 31st International Conference on Computational Linguistics (COLING 2025).
\item Abdulrahman Alalawi, Abdullah Alsuhaibani, Usman Naseem, Basem Suleiman, \textbf{Shoaib Jameel} and Imran Razzak. 2025. Few Labels with Active Learning: From Weak to Strong Labels for Misinformation Detection. The 1st International Workshop on  Large Language Models for Social Media (WebConf 2025)
\end{enumerate}

\section{UNDERGRAD RESEARCH GRANT} \textbf{Undergraduate Seed Research Grant:} My undergraduate university awarded me an INR (Indian Rupees) 10,000 seed grant (as PI) to support my future research. The grant could be used to buy research equipment or for conference travel. The award was prestigious and very competitive in the sense that I was the only student in my batch of 120 computer science students to be awarded this seed funding based on my undergraduate research accomplishments in the final year of my undergraduate study. Quality of the research proposal, followed by consistent academic performance, were also some of the criteria used to select the candidate.

\section{PhD GRANTS}
\begin{itemize}
    \item Scholarship name: BT and CSEE, University of Essex Doctoral Scholarship \\
    Duration: Three years (2019 -- 2022) \\
    Role: Main Supervisor and Principal Investigator
\end{itemize}

\section{INNOVATE UK RESEARCH GRANTS}
\begin{itemize}
    \item Proposal title: The utilisation of a novel combination of natural language processing (NLP) and image analysis to accelerate and automate the collection of Open Source Intelligence (OSINT), providing end-users with advanced insight into individuals without invading individual privacy.\\
Funding name: Knowledge Transfer Partnership (KTP)\\
Duration: Three years\\
Role: Lead Academic or Principal Investigator\\
Amount: \pounds350k \\
\textbf{Achievements:} My KTP Associate has already deployed a model on the company's production system last year.

    \item Proposal title: Develop AI methods to optimise interactions with customers. \\
    Funding name: Knowledge Transfer Partnership (KTP)\\
    Duration: Two years (2019 -- 2022). \textit{The project ended 10 months prematurely because the Associate found a permanent position at BT after demonstrating excellent performance.} \\
    Role: Support Academic or Co-Investigator (with 40\% role, more than the Principal Investigator) \\
    Amount: \pounds123k \\
    \textbf{Achievements:} The project resulted in: 1) model deployed in BT's internal production systems, 2) two research papers, 3) the Associate finding a permanent position at BT.

    \item Proposal title: The development of a new CPD tracker using AI and embedded machine learning to track and enhance the performance of all
staff. To gather data from the CPD tracker that will provide insight into effective strategies offering senior leaders a clear path for development driving industry-wide improvement. \\
    Funding name: Knowledge Transfer Partnership (KTP)\\
    Duration: Two years (2019 -- 2021) \\
    Role: Support Academic or Co-Investigator \\
    Amount: \pounds230k \\
    Note: \textit{This project was unfortunately withdrawn by the company midway as a result of financial losses incurred by it due to the COVID-19 pandemic.}
    
    \item Proposal title: The design and implementation of an integrated smart platform and bespoke client-side application for process integration (site plan material sharing, management and version control), enhancing effective office-site communications in the property development industry. \\
    Funding name: Knowledge Transfer Partnership (KTP)\\
    Duration: Two years (2020 -- 2022) \\
    Role: Support Academic or Co-Investigator (with 40\% role, more than the Principal Investigator) \\
    Amount: \pounds250k \\
    \textbf{Achievements:} A functional prototype is complete. The project won the \textcolor{magenta}{``Best KTP Project (SME) 2022''} award.
    
    \item Proposal title: Automating travel safety alerts. \\
    Funding name: Knowledge Transfer Partnership (KTP)\\
    Duration: Two years and six months (2021 -- 2024) \\
    Role: Academic Supervisor or Co-Principal Investigator (with 33.33\% role) \\
    Amount: \pounds230k \\
    Note: \textit{I wrote the complete work package and the KTP plan. While I was supposed to lead the project as PI, I decided to give other staff a chance due to my workload on other funded projects and allow a new face a chance to lead a KTP.}
\end{itemize}

\section{INNOVATION GRANTS}
\begin{itemize}
    \item Proposal title: \textbf{NG-CrAI: A Next-Gen Crowd-AI Framework for Scientific Discovery} \\
Funding name: \textbf{INESC TEC's Internal Seed Projects Grant} \\
Amount approved: \pounds35,938 \\
Duration: 2 years \\
Role: UK collaborator \\
Year: 2020--2021

    \item Proposal title: \textbf{Artificially Intelligent RoleCatcher, FinTex} \\
    Funding name: \textbf{Essex Innovation Voucher} \\
    Amount approved: \pounds10,000 \\
    Duration: One year \\
    Role: Principal Investigator \\
    Year: 2020--2021 \\
    \textbf{Achievements:} My work with Fintex Ltd., funded through the Essex Innovation Voucher, has been inspirational in various ways, e.g., how university and industry collaboration can lead to the successful building of a start-up, how novel models developed in research can be applied to industrial production systems, etc. The company, along with the university, has been covered by various media outlets discussing this success story. Fintex Ltd. is based at the Innovation Centre, University of Essex. This collaboration has been covered in various news media.
    
\item Proposal title: \textbf{Exploiting VR Technology and NLP for London Metropolitan Police} \\
Funding name: \textbf{London Metropolitan Police Seed Projects Grant} \\
Amount approved: \pounds50,000 \\
Duration: 1 year \\
Role: Co-Principal Investigator \\
Year: 2022

\item Proposal title: \textbf{DistillJob: Automatic Personalised Distillation of Job Search Results} \\
Funding name: \textbf{Web Science Institute Stimulus Fund 2022-23} \\
Amount approved: \pounds15000 \\
Duration: 6 months \\
Role: Principal Investigator \\
Year: 2023
\end{itemize}

\section{INDUSTRY FUNDING}
\begin{itemize}
\item Proposal title: \textbf{Automatic Geo-localization of Reading-for-Pleasure Children's Storybooks}\\
Funding name: \textbf{NVIDIA Academic Hardware Grants Program}\\
Hardware type: Quadro RTX6000 \\
Role: Principal Investigator

\item Proposal title: \textbf{Mining Time-Aware Urban Living Styles via Latent Semantic Concept Analysis}\\
Funding name: \textbf{Microsoft Research Asia Urban Informatics Research Fund, 2014}\\
Amount approved: USD 20,000\\
Duration: 1 year\\
Role: Co-Investigator \\
The content of this proposal was exclusively designed and written by me. I introduced a new dimension to the existing MSRA project called LifeSpec, which is available internally at MSRA. In particular, I added a temporal component when generating the urban human lifestyle patterns. To this end, I proposed an entirely new temporal model that improved upon the original urban human computing model proposed at MSRA. The proposed model improved upon the time and space complexities, quality of the urban patterns generated, handling the sparsity in the data and improving upon the quantitative tasks. I designed a non-parametric temporal graphical model and formulated the equations by myself. The two main challenges were technical formulations and the evaluation plans. I wrote the draft of the proposal with full technical details, the intuitions behind the model and how it could be useful to MSRA. Besides, I designed and prepared the slides to present the proposal to an audience from Microsoft Research, who came to judge all competing proposals. The committee applauded the idea and the model very much at the end of the presentation. The proposal cited my work, which was published in leading conferences, to convince the committee that we were competent enough to accomplish the required research task within a year.
\end{itemize}

\section{RESEARCH COUNCIL \& SMART GRANTS}
%I have contributed substantially towards two research grant proposals during my PhD study and postdoctoral work with my thesis advisor Prof. Wai Lam, who was the principal investigator in the projects. I was the investigator in these projects. Both proposals were accepted. The Microsoft Research grant was used solely to fund my postdoctoral research, and the RGC grant was used to support the research activities of my PhD research group.
\begin{itemize}
\item Proposal title: \textbf{PICASO: Probabilistic Conceptual Spaces Sensemaking under Uncertainty}\\
Funding name: Alan Turing Institute\\
Amount approved: GBP 140,000 \\
Completion date: Ongoing \\
Role: Principal Investigator

\item Proposal title: \textbf{Social-Aware News Trend Discovery via Joint Detection of Latent Information Structure from News
and Social Media Text Content}\\
Funding name: \textbf{HKRGC General Research Fund, 2014}\\
Amount approved: HKD 692,894 \\
Completion date: 31-12-2017 \\
Role: Co-Investigator \\
The RGC funding is among the few large research aids given by the Hong Kong government to some of the best research proposals. In this proposal, I handled and crafted the complete draft of the model's technicalities. The central theme of the proposal was to learn hierarchical human living and mobility patterns using data obtained from online social networks and apply the trained model to conduct event modelling and group recommendation. To this end, I designed the graphical model, the generative process and the mathematical formulations. I wrote the draft, which included the technical details, related work, the evaluation plan and the intuitions behind the model to convince the reviewers that the task and the proposed models are noteworthy and important enough to be funded.

\item Proposal title: \textbf{Multi-Modal Image Style Transfer: Automatically Geo-Localising Reading Material Digital Artwork for Increased Reader Engagement} \\
Funding name: Global Challenges Research Fund (GCRF) \\
Amount approved: \pounds36,5000 \\
Role: Principal Investigator \\
Start date: January 2021 \\
End date: July 2021

\item Proposal title: \textbf{AI modelled geo-localised and personalised reading-for-pleasure books in the UK schools} \\
Funding name: Royal Society Pairing Scheme 2022 \\
Role: Principal Investigator \\
News article\footnote{https://www.essex.ac.uk/news/2022/04/22/scientist-selected-for-royal-society-pairing-scheme}

\item Proposal title: \textbf{Research a fully automated employment law and industrial relations service enabled by Artificial Intelligence} \\
Funding name: Innovate UK Smart Grant 2022 \\
Amount approved: \pounds400k \\
Role: Principal Investigator \\
Start date: September 2022 \\
End date: March 2024
\end{itemize}

\section{STUDENT SUPPORT GRANTS}
\begin{itemize}
    \item Proposal title: \textbf{Automating Additional Reading Material Generation for Courses} \\
    Funding name: \textbf{Frontrunners Spring 2021} \\
    Amount approved: Undergraduate Summer Salary \\
    Role: PI
\end{itemize}

\section{GLOBAL RESEARCH ACHIEVEMENT}
\begin{itemize}
    \item One of the teams selected, \textbf{through a global competition}, for the ``Re-imagining Trust and Safety (T\&S) for Artificial Intelligence (AI)'' at United Nations Development Programme (UNDP)\footnote{https://www.undp.org/digital/ai-trust-and-safety}. The team is led by my PhD student, Ubaid Azam and me.
\end{itemize}


\section{PhD THESIS EXAMINATIONS}
\begin{itemize}
    \item Examined two PhD theses at the University of Essex
\end{itemize}
\begin{itemize}
    \item External examiner - committee also included Professor Udo Kruschwitz, University of Essex (\textit{currently Chair in Information Science at the University of Regensburg)}
    \begin{enumerate}
        \item Student: Andrius Mudinas, Birkbeck, University of London
        \item Year: 2019
        \item Supervisors: Dr. Dell Zhang and Professor Mark Levene 
    \end{enumerate}
    \item External Thesis Examiner - Sikkim Manipal University, India
    \begin{enumerate}
        \item Student: Papri Ghosh
        \item Year: 2019
        \item Supervisors: Dr. Tejbanta Chingtham and Professor MK Ghose
    \end{enumerate}
\end{itemize}

\section{INVITED TALKS}
\begin{itemize}
\item 2025 -- University of Sheffield Panel Talk
\item 2023 -- Thomson Reuters London Research Labs
\item 2023 -- Thomson Reuters India Research Labs [\href{https://medium.com/tr-labs-ml-engineering-blog/a-talk-on-contextualized-attention-embeddings-by-professor-shoaib-jameel-5e68e3bb12f1}{Link}]
\item AACL 2022 Birds-of-a-feather (BoF) Panel Discussion Member
\item DavisAI Seminar Talk
\item Invited talk on my GCRF project at CURE: Global Challenges (GCRF)
\item Invited talk at the 2021 BT AI Global Festival (Online event)
\item Lectures in the Summer School at Peking University, Beijing (July -- 2019). I was recently invited by the Department of Intelligent Computing and Sensing Laboratory at Peking University to deliver lectures in their summer school in Beijing, China. Nominations were first sent by the School of Computing, the University of Kent, to Peking University, which included my name as the only junior academic staff, followed by the names of the senior academic members, i.e., senior lecturers to professors. Academics from other UK and EU institutions had also sent their nominations. In the end, I was the \textbf{only academic} selected to deliver the lectures from the UK/EU region
\item Digital Humanities Workshop, CHASE Group (October -- 2018)
\item BenevolentAI (October -- 2018)
\item Centre for Language and Linguistics, the University of Kent (March -- 2018). I delivered the talk along with Professor Alessandro Vinciarelli from the University of Glasgow
\item Department of Computer Science and Information Systems, Birkbeck, the University of London (March -- 2018)
\item School of Mathematics, Cardiff University (November -- 2017)
\item Department of Systems Engineering and Engineering Management, The Chinese University of Hong Kong, SEEM Seminar Series (December -- 2017)
\item Workshop talk: ``Probabilistic Topic Models'' at the School of Computer Science and Informatics, Cardiff University (June -- 2017)
\end{itemize}

\section{IN THE MEDIA}
\begin{itemize}
    \item AI Music Lyrics Generation \href{https://www.thesun.co.uk/tech/23184161/ai-creates-first-love-song/}{[Link 1]}, \href{https://www.southampton.ac.uk/news/2023/07/love-letters-with-ai.page}{[Link 2]}
    \item Mining COVID-19 depression dynamics \href{https://www.gazette-news.co.uk/news/19068066.twitter-research-university-essex-shows-increase-depression-pandemic/}{[Link]}
    \item Essex Innovation Voucher collaboration with Fintex Ltd. \href{https://www.essex.ac.uk/news/2021/06/25/rolecatcher-helping-our-students}{[Link 1]}.
    \item Horus KTP project award \href{https://www.essex.ac.uk/news/2021/11/30/horus-security-consultancy-partners-with-university-of-essex}{[Link]}.
    \item My work on information retrieval as an undergraduate student \href{https://bashthebuilder.github.io/files/Shoaib.jpg}{[Link]}
    \item Our work with eBay led to the best paper nomination at CIKM-2021 \href{https://www.bandt.com.au/reaching-gen-z-a-new-tool-for-marketers/}{[Link]}. The project was selected by the eBay leadership team as a Top 3 Leaders' Choice Award recipient out of 2000 international submissions.
    \item I have published a write-up on Twitter's bias towards right-wing \href{https://theconversation.com/twitters-algorithm-favours-the-political-right-a-recent-study-finds-175154}{Link}.
\end{itemize}

\section{COMPETITIVE PROGRAMMING}
\begin{itemize}
    \item Initiated a ``unique'' student society at the University of Southampton called \textbf{Special Interest Group in Competitive Programming (SIGCompete)}.
    \begin{itemize}
        \item Trained students in competitive programming, such as ICPC, UKIEPC, NWERC and others throughout the university.
        \item Allocated special funding from the Head of School, ECS.
        \item Lead student teams abroad for competitive programming events.
    \end{itemize}
    \item \textcolor{blue}{Achievement:} Nominated as a \textbf{Jury} member in competitive programming; I formulate competitive programming questions for different events.
    \item \textcolor{blue}{Achievement:} Finalist (\textbf{one among the top 3}) in the prestigious \textbf{2024 Vice Chancellor's Awards} in the Teaching category (\textit{out of 630 applications}).
\end{itemize}

\section{REVIEWER SERVICE}
\textbf{Conferences}
\begin{itemize} \itemsep -10pt
    \item Senior Programme Committee Member
    \begin{itemize}
        \item International Conference on Acoustics, Speech, and Signal Processing (ICASSP -- 2025)
        \item International Joint Conference on Neural Networks (IJCNN--2025, IJCNN--2025)
        \item ACL Rolling Review (ARR) 2025, 2026
        \item European Conference on Artificial Intelligence (ECAI--2023, ECAI--2024, ECAI--2025)
        \item 38th AAAI Conference on Artificial Intelligence (AAAI--24, AAAI--25, AAAI--26)
    \end{itemize}
\end{itemize}
\begin{itemize} \itemsep -10pt
\item Programme Committee Member:
\begin{itemize}
\item ACS/IEEE International Conference on Computer Systems and Applications (AICCSA--2014)
\item Hawaii International Conference on System Sciences (HICSS--2015)
\item Russian Summer School in Information Retrieval (RussIR--2016)
\item World Wide Web Conference (WWW--2017, WWW--2018, WebConf--2024)
\item Knowledge Discovery and Data Mining (SIGKDD--2021, SIGKDD--2022, SIGKDD--2023, SIGKDD -- 2024)
\item The North American Chapter of the Association for Computational Linguistics (NAACL--2021)
\item Special Interest Group on Information Retrieval (SIGIR--2017, SIGIR--2018, SIGIR--2019, SIGIR--2020, SIGIR--2021, SIGIR--2022, SIGIR--2023, SIGIR--2024)
\item International Joint Conference on Artificial Intelligence (IJCAI--2017, IJCAI--2018, IJCAI--2019, IJCAI--2020, IJCAI--2023, IJCAI--2024)
\item ACM International Conference on Information and Knowledge Management (CIKM--2017, CIKM--2018, CIKM--2020, CIKM--2021, CIKM--2023, CIKM--2024)
\item Annual UK Workshop on Computational Intelligence  (UKCI--2017)
\item International Conference on Computational Linguistics (COLING--2018, COLING--2020, LREC-COLING--2024)
\item AAAI Conference on Artificial Intelligence (AAAI--2019, AAAI--2020, AAAI--2021)
\item European Conference on Information Retrieval (ECIR--2018, ECIR--2019, ECIR--2020, ECIR--2024)
\item Conference on Empirical Methods in Natural Language Processing (EMNLP--IJCNLP 2019, EMNLP--2020, EMNLP--2021, EMNLP--2022, EMNLP--2023)
\item Asia-Pacific chapter of the Association for Computational Linguistics (AACL--IJCNLP 2020, 2023)
\item Web Search and Data Mining (WSDM--2021, WSDM--2022, WSDM--2024)
\item European Chapter of the Association for Computational Linguistics (EACL--2021)
\item Annual Meeting of the Association for Computational Linguistics (ACL--2019, ACL--2020, ACL--2021, ACL--2023)
\item NeurIPS 2022, 2023
\item International Conference on Machine Learning (ICML--2023, ICML--2024)
\item European Conference on Machine Learning and Principles and Practice of Knowledge Discovery in Databases (ECML-PKDD--2023)
\item SIAM International Conference on Data Mining (SDM--2024)
\item SIGCHI 2024
\item The International Conference on Learning Representations (ICLR--2021, ICLR--2022, ICLR--2024) \\
\end{itemize}
\item External (Sub) Reviewer
\begin{itemize}\itemsep -10pt
\item Conference on Artificial Intelligence (AAAI--2017 and AAAI--2016)\\
\item International Joint Conference on Artificial Intelligence (IJCAI--2016 and IJCAI--2015)\\
\item Conference on Information and Knowledge Management (CIKM--2016)\\
\item World Wide Web Conference (WWW--2015 and WWW--2014)\\
\item Pacific-Asia Conference on Knowledge Discovery and Data Mining (PAKDD--2015)\\
\item Special Interest Group on Information Retrieval (SIGIR--2015, SIGIR--2014, SIGIR--2013 and SIGIR--2012)\\
\item Association for Computational Linguistics (ACL--2014)\\
\item Knowledge Discovery and Data Mining (KDD--2014 and KDD--2013)\\
\item IEEE International Conference on Data Mining series (ICDM--2013)
\end{itemize}
\end{itemize}
\textbf{Journals}
\begin{itemize} \itemsep -10pt
\item Transactions on the Web (TWEB) Distinguished Review Board member \\
\item WIREs Data Mining and Knowledge Discovery - \textcolor{red}{\textit{Awarded a certificate highlighting my reviewing contributions}} \\
\item Entropy \\
\item Fuzzy Sets and Systems \\
\item Transactions on Information Systems (TOIS) \\
\item ISPRS International Journal of Geo-Information \\
\item Information Retrieval Journal \\
\item The Computer Journal \\
\item Computer Speech \& Language - \textcolor{red}{\textit{Outstanding Reviewer recognition}}\\
\item Journal of Experimental \& Theoretical Artificial Intelligence\\
\item Information Processing and Management (IP and M)\\
\item Knowledge-Based Systems (KBS) - Elsevier \\
\item Neural Networks - Elsevier \\
\item Transactions on Asian Language Information Processing (TALIP) \\
\item Transactions on Knowledge and Data Engineering (TKDE) \\
\item Information (\textit{Member of the Reviewer Board}) \\
\item Computational Linguistics (\textit{Member of the Reviewer Board}) \\
\item IEEE Transactions on Computational Social Systems \\
\item Artificial Intelligence Communications
\end{itemize}

\section{JOURNAL EDITORIAL}
\textbf{Associate Editor}: AI Communications

\textbf{Guest Editor}: IEEE Transaction on Computational Social Systems\\
Topic: Special Issue on Computational Social Systems for COVID-19 Emergency Management and Beyonds \\
Year: 2020

\textbf{Guest Editor}: Frontiers in Research Metrics and Analytics \\
Topic: Multi-scale Assessment of Research Priorities over Acknowledgement and Funding Data during the COVID-19 Pandemic \\
Year 2022

\section{AWARDS AND ACHIEVEMENTS}
\textbf{University of Southampton}
\begin{itemize} \itemsep -10pt
\item Finalist -- Vice-Chancellor's Awards (Teaching)
\end{itemize}
\textbf{University of Essex}
\begin{itemize} \itemsep -10pt
\item Finalist in the ``ACS Digital Disruptors Awards 2022'' for our work on ``Onboarding Gen Z to eBay using an innovative multi-modal hypergraph model'' \\
\item Outstanding Early-Career Researcher Award winner in the Faculty of Science and Health -- Runner-up. \textit{Award Criteria:} Quality and quantity of research outputs, e.g., publications and attracting research income
\end{itemize}
\textbf{The Chinese University of Hong Kong}
\begin{itemize} \itemsep -10pt
\item Finalist in the ACM-Hong Kong Best Research Awards-2013 held at Hong Kong University of Science and Technology, Hong Kong. \textit{Only two best research candidates from the entire Faculty of Engineering of CUHK were selected to represent the university in the final}. The event was organised by ACM (Hong Kong Chapter) and Microsoft Research Asia\\
\item SIGIR-2013 Student Travel Award\\
\item Student Volunteer at SIGIR-2013 (received registration fee waiver) \\
\item ECIR-2013 Student Accommodation Grant offered by Yandex, Moscow\\
\item Best paper award at Beijing-Hong Kong International Doctoral Forum-2011\\
\item Department of Systems Engineering and Engineering Management Postgraduate Scholarship - 2009
\end{itemize}

\textbf{Undergraduate and before}
\begin{itemize} \itemsep -10pt
\item Letter of Appreciation for excellent performance at Tata Steel R\&D (India) - 2009\\

\item Nominated for the prestigious Overseas Research Students Awards Scheme (ORSAS), University of London (\textit{declined}) - 2009\\

\item Featured in the Indian national news for my ThinkSearch algorithm - 2008. I was interviewed one-on-one by a journalist from a leading Indian national newspaper, ``Hindustan Times''. The interview was published on the front page of my city's newspaper\\

\item Scored a GPA of 10/10 in the final semester during the undergraduate study - 2008\\

\item Best paper award at TechFest, Sikkim Manipal Institute of Technology - 2006\\

\item Third prize in SUPW/Science Workshop, Kerala Samajam Model School - 2002\\

\item Second prize in Rotary Club of Jamshedpur Computer Quiz - 2001\\

\item Consolation prize in SUPW/Science Workshop, Kerala Samajam Model School - 2001\\

\item Awarded at the Annual Academic Prize Night for scoring 90\% and above in the respective courses in ICSE (Indian Certificate of Secondary Education - Year 10) examination - 2001\\

\item Commendation certificates for an excellent academic record in each terminal examination - 1996-2000\\

\item Ranked among top three in the entire school - 1996-2000\\

\item Awarded at the Annual Academic Prize Night for an excellent academic track record in the entire year consistently from 1996 till 2000\\

\item Third prize in Inter-School Academic Quiz, Kerala Samajam Model School - 2000\\

\item First prize in Science Quiz, Kerala Samajam Model School - 1999\\

\item Featured in the national news for securing the top rank in school - 1996. The newspaper carried a lengthy story about me when I stood first in the entire batch for the first time. What was so inspiring about the story was that I was just a below-average student a year back\\

\item Best student in Discipline in the entire year in Taekwondo - 1995 \\

\item Rewarded numerous times in Taekwondo championships, belt examination and camps
\end{itemize}

\section{TEACHING EXPERIENCE}
\textbf{University of Southampton}\\
\begin{enumerate}
    \item COMP3222 and COMP6246 -- Machine Learning Technologies (Module Lead). \textbf{Achievements:} I have made significant contributions ranging from creating additional pre-recorded lecture materials and a Q\&A document for each topic to help students prepare for the examination and additional lecture notes. The module has consistently received positive feedback.
    \item COM6256 -- Knowledge Graphs for AI Systems (Module Lead). \textbf{Achievements:} The module has been consistently ranked among the top 1\% of modules in the entire school. This is because I have demonstrated excellent dedication towards providing student support.
    \item GENG0015 -- Engineering Foundations (Teaching Support).
\end{enumerate}

\textbf{University of Essex, Colchester}\\
\begin{enumerate}
    \item CE807 - Text Analytics (Module Lead). \textbf{Achievements:} This module received excellent student feedback and a teaching score. Student comments and scores can be found \href{https://drive.google.com/file/d/1S5xoQOVRjkppus1jlcfoXiVYurx1Uhpj/view?usp=sharing}{[here]}.
    \item CE306/CE706 - Information Retrieval (Second Module Lead)
    \item CE299 (Project Module) - Lab Supervisor
\end{enumerate}
\textbf{University of Kent, Medway}\\
\begin{enumerate}
    \item CO334 - People and Computing (Module Convener)
    \item CO547 - Agile Software Development and Security (Module Convener)
    \item CO639 - Electronic Commerce
    \item CO644 - Semantic Web
\end{enumerate}
\textbf{Students' feedback about my teaching at the University of Kent:} \newline
Immediately sent by a student after my first lecture ever in Kent (Module CO639) to his tutor who is also a lecturer in the School of Computing at Medway: \say{\textit{Afternoon mrs, I would just like to put in a small word of gratitude \& say today's lecture felt very engaging, not saying that yours or XXX (other lecturer's name made confidential) aren't but, considering it was his (didn't get his name) first time lecturing here, he engaged well I believe - he ensured we understood what was briefed in the powerpoints and just made things easier. Perhaps it was today's topic, that made the teaching easier, I don't know but I just wanted you to pass my regards to him, I just felt more engaged than usual today... thank you}}.

This is what the tutor of the student had to say: \say{\textit{Congratulations Shoaib! This feedback is from a student who never ever contacts me usually (and he’s my tutee!) and has not really been engaging with several courses. We need to observe your teaching rather than the other way around, it seems :)}}

Another student from a different module that I taught next semester (Module CO334). It shows that I have maintained a good teaching track record:

\say{\textit{Hi, Shoaib, I am contacting you regarding my review of your teaching style. I really enjoy attending your lecture and prefer this approach to teaching, however, I feel sometimes it can go off-topic. I enjoy sharing my opinion and having the class getting more involved gives the lecture a better atmosphere. Furthermore, I think using videos makes it easier to write notes and is more informative than having 80+ slides to go through. I would love to see examples incorporated into this video to showcase what you're saying rather than having to see you draw. Thanks.}}

Another student from the same module as above:

\say{\textit{Dear Shoaib,
Just wanted to say that as part of your People and Computing lecture on a Friday morning, I really do enjoy your teaching. You seem more charismatic than almost all of the other lecturers and I think the way you teach is engaging too. The only feedback i can give is not about you but is infact about the module being a little basic. I understand that you cannot change the stuff that has to be taught so maybe try and engage the class by creating interactive activities. I'm not too sure what other people in the class would think about it but I know it would help me become more awake as it is a Friday morning after all!
Hope that can be of some help,}}

This is what one mature student had to say about my teaching in CO639:

\say{\textit{Hi Shoaib, 
I wanted to let you know that I really enjoyed the assignment on cryptocurrencies. I really like your style of teaching. The idea of a discussion first gets you thinking and the slides afterwards cement it in your mind.
The only negative I would say is you ask too much if your teaching style is okay.  I totally understand why you do this but it can be a little distracting. I think it would be better if you just rolled with it and if there was an issue I'm sure the students would let you know.
I like you as a person and a lecturer and wish you all the best in the future.}}

The comment below shows that I have consistently maintained my innovative teaching methods in my second year as Lecturer. A final year student from the module Semantic Web that I am teaching in 2019.

\say{\textit{Hi Shoaib, 
Just wanted to say I loved the lecture today! And your methods have helped us a lot. Thank you! }}

\textbf{The Chinese University of Hong Kong}\\
Introduction to Engineering - I (ENGG1100)
\begin{itemize}\itemsep -3pt
\item Supervised students in designing and constructing a robotic car for automatic navigation
\item Taught students about the basics of 3D printing and using SolidWorks
\end{itemize}
Information Systems Management (SEEM 3490)
\begin{itemize}\itemsep -10pt
\item Delivered lectures on Information Security and Cryptography \\
\item Conducted revision classes covering the entire syllabus before the final examination (\textit{the most challenging part!}) \\
\item Supervised student projects
\end{itemize}
Data Structures (CSC2100E/F)
\begin{itemize}\itemsep -10pt
\item Delivered lectures on Computational Complexity, Pseudo-code writing, Searching and Sorting \\
\item Formulated interesting assignment questions based on real-life problems\\
\item Conducted revision classes covering the entire syllabus before the final examination (\textit{again the most challenging part!})\\
\item Supervised student programming projects
\end{itemize}
\textbf{Sikkim Manipal Institute of Technology}\\
Data Structures - Functions and Pointers in C
\begin{itemize}\itemsep -10pt
\item Organized student workshops covering all aspects of ``Functions'' and ``Pointers'' in the C programming language
\end{itemize}

\section{GRADUATED PhD SUPERVISION}
\begin{itemize}
\item \textbf{Ph.D. student}: Liao Yi (CUHK) - I can be best described as a ``technical advisor'' to them. My responsibility was to advise them on mathematical modelling, scientific programming and experiment design with various analyses. The student is currently working at Huawei's Research Lab as a Research Scientist
\item \textbf{Masters (MPhil)}: Sam Lai (CUHK) - I have supervised Sam for a year and have taught him about data analysis and crawling for online social networks. I have also trained him in probabilistic topic modelling programming in Java. He has successfully defended his MPhil thesis, and subsequently, he will take the role of a research assistant in my PhD thesis supervisor's group. We have jointly published a paper in a flagship information retrieval conference recently (SIGIR).
\item \textbf{Ph.D. student}: Mahsa Abazari Kia, School of Computer Science and Electronic Engineering, The University of Essex. (\textit{Research Supervisor}). \textbf{Accomplishments:} 1) Model development complete to be deployed on BT systems; 2) Lecturer at a university in London.
\item \textbf{Ph.D. student}: Reyhaneh Hashempour, School of Computer Science and Electronic Engineering, The University of Essex. (\textit{Research Co-Supervisor}). Presently working at BT as a Data Scientist.
\item \textbf{PhD student}: Ant\'onio Correia, Ph.D. candidate in Computer Science, University of Tras-os-Montes e Alto Douro, Vila Real, Portugal
Researcher at INESC TEC, Porto, Portugal. (\textit{Research Co-Supervisor}). \textbf{Accomplishments:} 1) PARSUK Xperience scholarship holder; 2) FCT (Portuguese Foundation for Science and Technology) international research visit travel grant winner; 3) Scientific Production Award 2021; 4) FLAD Science Award USA Grant; 5) Portuguese Foundation for Science and Technology’s Award of Research Grants for Doctoral Degree; 6) Research Grant Holder under the projects ``Video Interactivo de Comunicacao - VIC''; 7) University of Kent's Postgraduate Research Internship; 8) Portuguese Minister of Economy, Innovation and Development's IAPMEI funded projects ``SciCrowdML: A Crowd-Computing, Intelligent Big Data Analysis System for Students and Scholars''; 9) Best Paper Award (1st place) at the 4th International Research Conference on Virtual Worlds; 10) \textbf{First} Portuguese PhD student to be awarded the prestigious \textcolor{magenta}{Microsoft Research PhD Fellowship 2021}; 11) Top 10 in ``The Annual Excellence in Doctoral Research Award'' [\href{https://www.youtube.com/watch?v=9Wvq0P95560}{Video link}] by  iSCSi - International Conference on Industry Sciences and Computer Sciences Innovation.
\item \textbf{PhD student}: Hamad Zogan, PhD candidate in School of Computer Science \& Advanced Analytics Institute, Faculty of Engineering \& IT, University of Technology Sydney. Australia. (\textit{Research Co-Supervisor}). \textbf{Accomplishments:} 1) Work appeared in the various news media; 2) A full paper at prestigious SIGIR-2021 conference

\item \textbf{PhD student}: Mozhgan Talebpour, PhD candidate in the School of Computer Science and Electronic Engineering, The University of Essex. (\textit{Research Supervisor}). 
\textbf{Accomplishments:} Computational model deployed on the company's production systems, significantly aiding in its everyday business. Research papers published in premier conferences (ECIR and SIGIR)
\end{itemize}

\section{CURRENT RESEARCH SUPERVISION}
\begin{itemize}
\item \textbf{PhD student}: Lingzhi Shen. (\textit{Research Supervisor}).
\item \textbf{PhD student}: Ubaid Azam. \textcolor{magenta}{Winner: Doctoral College Director's Award 2026 -- Public Engagement \& Outreach and Knowledge Exchange \& Enterprise} (\textit{Research Supervisor}).
\item \textbf{PhD student}: Inamullah Khan. (\textit{Research Supervisor}).
\item \textbf{PhD student}: Bador Ali Al Sari. (\textit{Research Supervisor}).
\item \textbf{PhD student}: Dalia Nahhas. (\textit{Research Supervisor}).
\item \textbf{PhD student}: Ibtisam Khader Alzahrani. (\textit{Research Supervisor}).
\item \textbf{PhD student}: Fatima Alshareef. (\textit{Research Supervisor}).
\item \textbf{PhD student}: Ghadir Alselwi. (\textit{Research Co-Supervisor}).
\item \textbf{PhD student}: Abdulrahman Alalawi. (\textit{Research Co-Supervisor}).
\item \textbf{PhD student}: Abdulsalam Obaid A. Alharbi (\textit{Research Co-Supervisor}).
\item \textbf{PhD student}: Matthew Javanshir. (\textit{Research Co-Supervisor}).
\end{itemize}

\section{COMPLETED RESEARCH SUPERVISION}
\begin{itemize}
    \item \textbf{Postdoctoral KTP Associate}: Yogesh Kumar Meena. School of Computer Science and Electronic Engineering, The University of Essex. Currently, Assistant Professor at the Indian Institute of Technology in Gujarat.
    \item \textbf{KTP Associate}: Faizan Khan. School of Computer Science and Electronic Engineering, University of Essex. KTP with Horus Security Consultancy in Oxford. Currently working for 3 UK as a Data Scientist.
    \item \textbf{Research Assistant}: Anastasia Tomita. School of Computer Science and Electronic Engineering, The University of Essex. Currently, a Graduate Trainee at Lloyds Bank.
    \item \textbf{Research Assistant}: Yaseen Mohammed Osman. Electronics and Computer Science, University of Southampton. Currently, a PhD student at the University of Southampton.
\end{itemize}

\section{COLLEGIALITY}
\begin{itemize}
    \item \textbf{CO320: Introduction to Object-Oriented Programming} Took over Java tutorial classes for the undergraduate students; covered on behalf of a staff member who abruptly resigned, leaving no one to cover the required tutorials. (the University of Kent in 2019)
    \item \textbf{CO334: People and Computing} Took over the lectures for the staff who was off work due to paternity reasons. (the University of Kent in 2018)
    \item \textbf{CO600: Final-year Undergraduate Major Projects} Supervised a group because the main supervisor went on sabbatical leave. (the University of Kent in 2018)
    \item \textbf{Admissions Officer} Took over the main undergraduate responsibility because the admissions officer was offered the Head of School role. (the University of Kent in 2018)
    \item \textbf{New Lecturer interviews} Took the responsibility on very short notice of taking care of the applicants who visited the campus for the new Lecturer interviews. (the University of Kent in 2018)
    \item \textbf{Examination script marking} Marked final-year written undergraduate examination answer scripts on behalf of an academic staff who suddenly took leave. (the University of Kent in 2018 and 2019)
    \item \textbf{Data Science group logo design} Took responsibility for designing the group's logo. (the University of Kent in 2018)
\end{itemize}

\section{PROFESSIONAL EXPERIENCE}
\textbf{Tata Steel Limited}, Jamshedpur, India \\
Intern at the Research and Development Division (R\&D) \hfill {\textit{January 2008 - June 2009}}
\begin{itemize}\itemsep -10pt
\item Designed, developed and deployed the CPU load balancing architecture for the Linux cluster\\
\item Designed, developed and deployed a multi-threaded local Intranet crawler\\
\item Designed and developed the ``document-streaming'' algorithm\\
\item Estimation of yield-point from hot-strip mill data\\
\item Drafted a detailed technical report on text classifiers
\end{itemize}
\textbf{Tata Steel Limited}, Jamshedpur, India \\
Intern at the Automation Division \hfill {\textit{July 2007 - August 2007}}
\begin{itemize}\itemsep -10pt
\item Used Data Mining to come up with the relationship between steel casting parameters and the billet caster
\end{itemize}

\section{GRADUATE COURSES UNDERTAKEN}
\begin{itemize}\itemsep -10pt
\item Techniques for Data Mining (CSC 5180)\\
\item Text Mining Models \& Applications (SEG 5680)\\
\item Advanced Database \& Information Systems (SEG 5010)\\
\item Knowledge Systems (SEEM 5470 )\\
\item Logistics \& Transport Planning (SEEM 5600)
\end{itemize}

\section{PROFESSIONAL TRAINING}
\begin{itemize}\itemsep -10pt
\item Postgraduate Certificate for Higher Education programme (PGCHE) - \textbf{Status: Awarded (\textit{Fellow of the HEA}).} Two compulsory and two optional modules were completed, thereby fulfilling the requirements for the PGCHE. \textit{Out of four modules, DISTINCTION in two.} \\
\item Grant Writing Workshops \\
\item Unconscious Bias \\
\item Engaging with UK Parliament and Government \\
\item Communication Skills Training\\
\item Teaching Assistant Training
\end{itemize}

\section{SERVICES AND PARTICIPATION}
\textbf{University of Southampton}
\begin{itemize}\itemsep -10pt
    \item Deputy Appeals Officer \\
    \item Founder of Special Interest Group in Competitive Programming (SIGCOMPETE)\\
    \item PhD progress committee panel member
\end{itemize}
\textbf{University of Essex}
\begin{itemize}\itemsep -10pt
    \item Social Media Coordinator \\
    \item Degree Apprentice Coordinator \\
    \item Placement Year Coordinator \\
    \item Workload Allocation Model committee member \\
    \item Data Science and AI Project Placements: The Masters DAIV, project selection panel member \\
    \item AI for Decision Making group’s blog writer on Medium platform to publicise every group member's work \\
    \item Online teaching delivery task force group member \\
    \item Member of Benefits and Services Subcommittee, IEEE \\
    \item Essex Open Day lecture delivery \\
    \item MSc Course validation panel member
\end{itemize}
\textbf{University of Kent}
\begin{itemize}\itemsep -10pt
\item Mentored three students at ACL-2019 \\
\item Data Science group web page and blog maintainer\\
\item Undergraduate admissions work which includes participation in Open Days (usually Saturdays), Applicant Days, Undergraduate interviews, Ad-hoc undergraduate interviews, admission meetings, undergraduate campus tours, UCAS online applications processing, processing undergraduate student applicant day invitations - \textit{This is a major duty and the most important administration work in the school. I continued to maintain a strong teaching and research track record despite taking up this time-consuming role}\\
\item Participated in organising the new faculty recruitment \\
\item Session Chair at AAAI-2019 (Robotics Track) and SIGIR-2021 (Social Track), KDD-2021 (Web Mining Track) \\
\item Panel member in student plagiarism hearing\\
\item Helping the university with student campus recruitment\\
\item Panel member in Medway disciplinary committee
\end{itemize}
\textbf{Conference Organization Experience}
\begin{itemize}\itemsep -10pt
\item ``Board of members for Conference Programme editorial'', Thermec'2013, Las Vegas, Dec. 2 to 6, 2013\\
\item Organizing Chair, Beijing-Hong Kong International Doctoral Forum (2011)\\
\item Technical Paper Reviewer, Beijing-Hong Kong International Doctoral Forum (2010-2011)
\end{itemize}
\textbf{Other Services}
\begin{itemize}\itemsep -10pt
\item University Open Day planning, organization and execution (CUHK)\\
\item Student Volunteer at the 2014 Big Data Workshop jointly organized by KAIST-Microsoft Research Collaboration Center (KMCC), Korea; Department of Computer Science, Korea Advanced Institute of Science and Technology; the Key Lab of High Confidence Software Technologies (CUHK Sub-Lab), Ministry of Education, China; and the Department of Systems Engineering and Engineering Management, The Chinese University of Hong Kong (SEEM-CUHK) from April 22-23, 2014\\
\item Constituted a Special Interest Group (SIG) in information retrieval at SMIT\\
\item Regularly assisted the computer laboratory in charge of hardware and software maintenance at SMIT\\
\item Active member of the students' technical body, ACCESS at SMIT\\
\item Performed several social services such as teaching computers to backward people in India\\
\item Lifetime member of the Numismatics Club of Jamshedpur, India\\
\item An active participant in debates, spelling contests, etc
\end{itemize}


\section{PROGRAMMING SKILLS}
~~Programming and Scripting Languages
\begin{itemize}\itemsep -10pt
\item C, C++, MATLAB, R, Python\\
\item SED, AWK, PERL, Linux Shell Scripting
\end{itemize}

Productivity Applications
\begin{itemize}\itemsep -10pt
\item \LaTeXe{}, MS Office, Libreoffice, Vim
\end{itemize}

Operating Systems
\begin{itemize}\itemsep -10pt
\item Linux and its variants, Microsoft Windows, Sun Solaris
\end{itemize}

\end{resume} 
\end{document} 
